\chapter{Stoma Reversal}

\section*{Introduction \& Clinical Indications}
Stoma reversal is a reconstructive surgical procedure aimed at restoring normal gastrointestinal continuity following the creation of a temporary diverting ostomy, which may be a colostomy or an ileostomy. This operation is performed to eliminate the external stoma, allowing for the physiological passage of stool through the anus. The procedure involves mobilizing the stoma, resecting any non-viable or scarred bowel tissue, and rejoining the proximal and distal segments of the bowel. The primary clinical significance of stoma reversal lies in its ability to significantly enhance a patient's quality of life by alleviating the burdens associated with stoma management, such as appliance care, skin irritation, and body image concerns. Stoma reversal is typically indicated once the underlying condition that necessitated the initial ostomy has resolved, and the distal bowel segment is healthy and functional.

\begin{itemize}
  \item Colostomy closure
  \item Ileostomy closure
  \item Hartmann's procedure reversal
  \item Bowel continuity restoration
  \item Enterostomy takedown
  \item Stoma take-down
  \item Ostomy closure
\end{itemize}

The fundamental principle of stoma reversal is to re-establish a functional, tension-free, and well-vascularized connection between the proximal and distal segments of the gastrointestinal tract. This involves several critical steps, each guided by the rationale of minimizing complications and optimizing functional outcomes. 

Initially, the stoma, along with a small segment of the bowel wall that formed the stoma, is carefully mobilized from the abdominal wall. This mobilization must be meticulous to avoid injury to adjacent structures, such as the small bowel, ureters, or major vessels, and to ensure adequate length of healthy bowel for a tension-free anastomosis. Often, the stoma-bearing segment of the bowel is resected to remove any scarred, inflamed, or ischemic tissue, providing fresh, healthy margins for the subsequent connection, thereby reducing the risk of anastomotic breakdown.

Following mobilization and potential resection, the two healthy bowel ends are brought together and joined to create an anastomosis. This connection can be performed either by hand-sewn sutures or using surgical stapling devices, typically in an end-to-end, side-to-end, or end-to-side configuration, depending on the specific anatomy, bowel caliber discrepancy, and surgeon preference. A critical technical goal is to create a secure, patent, and well-vascularized anastomosis to minimize the risk of leakage, which is the most feared complication. Finally, the defect in the abdominal wall fascia, created by the original stoma, is meticulously repaired to prevent future incisional hernias, and the skin is closed. The ultimate aim is to restore the normal physiological pathway for defecation, allowing stool to pass through the colon and rectum to the anus, thereby significantly improving the patient's quality of life.

The concept of creating a stoma to divert fecal flow has ancient roots, with early descriptions dating back to the 18th century by surgeons like Amussat and Littr\'{e} for cases of intestinal obstruction. However, these were primarily permanent solutions, often associated with high morbidity and mortality. The modern era of temporary stomas and their subsequent reversal gained prominence in the 20th century, particularly with advancements in anesthesia, aseptic techniques, and understanding of fluid and electrolyte balance. 

Henri Hartmann's description of a procedure in 1921 for left-sided colonic obstruction or perforation, initially involving a permanent colostomy and closure of the distal stump, laid the groundwork for planned staged resections. Its utility as a temporary measure, followed by a planned reversal, became increasingly recognized in the latter half of the century as surgical techniques and patient care improved.

The widespread adoption of temporary diverting loop ileostomies to protect distal colorectal anastomoses, especially after low anterior resections for rectal cancer, significantly increased the frequency of stoma reversals. Pioneers in colorectal surgery, such as Sir Alan Parks and John Goligher, refined techniques for safe anastomosis and abdominal wall closure, emphasizing tension-free connections and adequate blood supply. The evolution from open laparotomy to minimally invasive laparoscopic approaches for both stoma creation and reversal has further improved patient outcomes, reducing pain, hospital stay, and recovery time. Laparoscopic reversal, first reported in the early 1990s, has become increasingly common, offering benefits such as smaller incisions, reduced adhesion formation, and faster return to normal activity. The understanding of factors influencing anastomotic healing and the importance of careful patient selection have been key drivers in the successful evolution of stoma reversal as a routine and safer reconstructive procedure.

Stoma reversal is indicated when the clinical conditions that necessitated the creation of a temporary stoma have been successfully resolved, making the stoma no longer necessary for patient safety or recovery. The decision to proceed with reversal is based on a comprehensive assessment of the patient's overall health and the status of the underlying pathology.

\begin{itemize}
  \item Resolution of acute inflammatory processes, such as severe diverticulitis, peritonitis, or inflammatory bowel disease flares, allowing for safe re-anastomosis without active inflammation.
  \item Healing and maturation of a previously created distal colorectal or ileal anastomosis, typically confirmed by imaging (e.g., Gastrografin enema) or endoscopic evaluation (e.g., flexible sigmoidoscopy) to rule out leak or stricture.
  \item Completion of adjuvant therapies, such as chemotherapy or radiation, following resection for colorectal malignancy, where the stoma was protective during the treatment phase.
  \item Resolution of peritonitis, intra-abdominal sepsis, or severe trauma that initially required fecal diversion, ensuring a clean and stable abdominal environment.
  \item Significant improvement in the patient's overall nutritional status and general health, optimizing conditions for surgical healing and reducing perioperative risks.
  \item Patient desire for improved quality of life, freedom from stoma appliance management, and restoration of normal body image, provided clinical criteria are met.
  \item Closure of a loop ileostomy created after restorative proctocolectomy for ulcerative colitis or familial adenomatous polyposis, once the ileal pouch-anal anastomosis has healed and matured.
  \item Resolution of distal bowel obstruction that was initially managed with a diverting ostomy, confirming patency of the distal segment.
\end{itemize}

Stoma reversal is unequivocally a secondary, reconstructive surgical procedure. It is never performed as a primary treatment for an acute disease or condition. Instead, it is undertaken only after an initial, definitive surgical intervention or the resolution of an acute medical crisis that necessitated the creation of a temporary diverting ostomy. 

For instance, a patient might undergo a primary anterior resection for rectal cancer, followed by a temporary loop ileostomy to protect the low anastomosis from potential leakage during the critical healing phase. The stoma reversal, in this context, is a subsequent, planned operation to restore the anatomy and function that were temporarily altered.

Similarly, a Hartmann's procedure for complicated diverticulitis with perforation and peritonitis is a primary, often life-saving, operation involving resection of the diseased colon and creation of an end colostomy. The stoma reversal is then a planned secondary procedure, typically performed several months later, once the patient has fully recovered, inflammation has subsided, and the distal rectal stump is deemed healthy. Its role is to restore normal physiological function and significantly improve the patient's long-term quality of life by eliminating the need for an ostomy appliance, rather than to treat the initial pathology. The timing of reversal is crucial and depends on the patient's overall health, the healing of the primary anastomosis, and the resolution of any underlying disease processes.

\section*{Relevant Surgical Anatomy}
A comprehensive understanding of the relevant surgical anatomy is crucial for the successful execution of stoma reversal. The procedure primarily involves the small intestine (ileum) or large intestine (colon), depending on the type of ostomy. 

For an ileostomy, the terminal ileum is involved, characterized by its relatively thin wall, abundant lymphoid tissue (Peyer's patches), and rich vascular supply from the superior mesenteric artery (SMA) via numerous arcades. The mesentery of the small bowel is extensive, allowing for significant mobilization. In the case of a colostomy, any segment of the colon can be involved, each with distinct anatomical considerations. The ascending and transverse colon receive blood supply from branches of the SMA (ileocolic, right colic, middle colic arteries), while the descending and sigmoid colon are supplied by the inferior mesenteric artery (IMA) via the left colic and sigmoid arteries. The marginal artery of Drummond, a crucial anastomotic arcade connecting the SMA and IMA territories, provides collateral circulation and must be preserved.

Venous drainage generally parallels the arterial supply, emptying into the superior and inferior mesenteric veins, which ultimately drain into the portal system. Lymphatic drainage follows the vascular bundles to regional lymph nodes. The nerve supply to the bowel is autonomic, with sympathetic innervation from the splanchnic nerves and parasympathetic from the vagus nerve (proximal colon) and pelvic splanchnic nerves (distal colon and rectum). 

The abdominal wall anatomy is also critical, comprising skin, subcutaneous fat, Scarpa's fascia, external oblique, internal oblique, and transversus abdominis muscles, their aponeuroses forming the rectus sheath, and the transversalis fascia and peritoneum. The previous stoma site will have a fascial defect that requires meticulous repair to prevent incisional hernia. Careful identification and preservation of the ureters, bladder, and major retroperitoneal vessels are essential, especially during mobilization of a distal colonic segment or in cases of extensive adhesions.

\section*{Preoperative Workup \& Preparation}
Thorough preoperative preparation is crucial for optimizing outcomes and minimizing complications associated with stoma reversal. This typically involves a comprehensive medical evaluation and specific interventions to ensure the patient is in the best possible condition for surgery.

\begin{itemize}
  \item \textbf{Medical Clearance:} A full history and physical examination, including detailed cardiac and pulmonary assessments, to ensure the patient is fit for general anesthesia and major surgery. This may involve an electrocardiogram (ECG), chest X-ray, pulmonary function tests, and consultation with an internist or cardiologist.
  
  \item \textbf{Laboratory Tests:} Routine blood work, including complete blood count (CBC), electrolyte panel, renal and liver function tests, coagulation profile (PT/INR, PTT), and blood type and screen, to assess overall health, identify any abnormalities, and prepare for potential transfusion.
  
  \item \textbf{Imaging and Endoscopy:} A contrast enema (e.g., Gastrografin enema) or flexible sigmoidoscopy is often performed to assess the patency, integrity, and health of the distal bowel segment, ruling out any strictures, fistulas, active inflammation, or recurrent malignancy.
  
  \item \textbf{Bowel Preparation:} Mechanical bowel preparation (e.g., polyethylene glycol solution) combined with oral antibiotics (e.g., neomycin and metronidazole) is frequently prescribed to reduce bacterial load in the colon, thereby decreasing the risk of anastomotic leak and surgical site infection.
  
  \item \textbf{Medication Review:} Adjustment or temporary discontinuation of certain medications, such as anticoagulants (e.g., warfarin, DOACs) or antiplatelet agents (e.g., aspirin, clopidogrel), is necessary to minimize bleeding risk. Diabetic medications may also need adjustment to optimize glycemic control.
  
  \item \textbf{NPO Status:} Patients are instructed to be NPO (nil per os) for a specified period (typically 6-8 hours) before surgery to prevent aspiration during anesthesia induction.
  
  \item \textbf{Prophylactic Antibiotics:} Intravenous broad-spectrum antibiotics are administered shortly before the incision to reduce the risk of surgical site infection.
  
  \item \textbf{Informed Consent:} Detailed discussion with the patient regarding the procedure, potential risks, benefits, alternatives, and expected postoperative course, ensuring full understanding and consent.
\end{itemize}

Stoma reversal, while generally safe, has specific contraindications and requires careful precautions to ensure patient safety and successful outcomes. The decision to proceed must weigh the benefits of reversal against the potential risks.

\textbf{Situations requiring optimization, delay, or an alternative plan:}
\begin{itemize}
  \item Uncontrolled sepsis or active intra-abdominal infection that could compromise anastomotic healing and lead to catastrophic leak.
  
  \item Severe malnutrition or immunosuppression may warrant optimization or a staged/alternative plan, but the decision depends on urgency and the reason for the stoma.
  
  \item Distal bowel obstruction or unresectable malignancy in the distal segment, which would prevent normal passage of stool and lead to recurrent obstruction.
  
  \item Active distal inflammatory bowel disease requires disease-specific assessment and may favor delay or another reconstruction; it is not a universal absolute contraindication.
  
  \item Severe, uncorrectable cardiac, pulmonary, or renal comorbidities that render the patient unfit for general anesthesia and major abdominal surgery.
  
  \item Inadequate length or viability of the proximal or distal bowel segments to create a tension-free, well-vascularized anastomosis.
\end{itemize}

\textbf{Relative Contraindications and Precautions:}
\begin{itemize}
  \item Previous pelvic radiation therapy, which can lead to poor tissue quality, impaired healing, and increased risk of anastomotic complications in the distal bowel.
  
  \item Extensive intra-abdominal adhesions from prior surgeries, increasing the risk of bowel injury, prolonged operative time, and difficulty achieving a tension-free anastomosis.
  
  \item Obesity (BMI $>$ 30-35 kg/m\textasciicircum2), which can complicate surgical access, increase operative time, and elevate the risk of wound complications, infection, and hernia.
  
  \item Chronic steroid use or other immunosuppressive medications, requiring careful management, potential temporary cessation, or stress-dose steroids to mitigate adrenal insufficiency.
  
  \item Poorly controlled diabetes, which can impair wound healing, increase infection risk, and contribute to anastomotic complications.
  
  \item Patient refusal or lack of understanding regarding the risks and benefits of the procedure, necessitating further counseling.
  
  \item Inadequate assessment of the distal bowel, necessitating further imaging or endoscopy before proceeding to ensure patency and health.
  
  \item Recurrent or persistent underlying disease that may necessitate further surgical intervention or make reversal unsafe.
\end{itemize}

\section*{Surgical Technique \& Procedure}
The procedure for stoma reversal is performed under general anesthesia, ensuring the patient is completely unconscious and pain-free throughout the operation. The patient is typically positioned supine, though a modified lithotomy position may be used if access to the anus is required for specific anastomotic techniques (e.g., transanal stapling) or leak testing. Prophylactic intravenous antibiotics are administered prior to incision.

The surgical approach usually begins with a circum-stomal incision, carefully dissecting around the stoma to mobilize the proximal and distal bowel segments from the abdominal wall and any surrounding adhesions. This mobilization must be meticulous to avoid injury to adjacent structures. In more complex cases, or if there is concern for extensive intra-abdominal adhesions, a formal midline laparotomy may be required to gain adequate exposure and safely lyse adhesions. 

Once mobilized, the stoma-bearing segment of the bowel, which often contains scarred, inflamed, or ischemic tissue, is resected to ensure healthy, well-vascularized bowel ends for the anastomosis. The two healthy bowel ends are then brought together for anastomosis. This can be achieved using a hand-sewn technique with absorbable sutures (single or double layer) or with a surgical stapling device (linear or circular stapler), creating an end-to-end, side-to-end, or end-to-side connection, depending on the specific anatomical configuration and surgeon preference.

After the anastomosis is completed, its integrity is often tested to detect potential leaks. This can involve injecting air or fluid into the bowel lumen while submerging the anastomosis in saline, or by performing a flexible sigmoidoscopy with air insufflation and direct visualization. Finally, the fascial defect in the abdominal wall, where the stoma was located, is meticulously repaired to prevent the development of an incisional hernia. This may involve primary suture repair using non-absorbable sutures or, in some cases, the use of prosthetic mesh, particularly for larger defects or in high-risk patients. The subcutaneous tissue and skin are then closed, often with absorbable sutures, and a sterile dressing is applied. Drains are rarely placed unless there is significant contamination, concern for fluid collection, or a high-risk anastomosis.

\section*{Complications \& Risks}
As with any surgical procedure, stoma reversal carries inherent risks, which can be broadly categorized into general surgical complications and those specific to bowel surgery and stoma closure. Patients must be thoroughly informed of these potential adverse events.

\textbf{General Surgical Risks:}
\begin{itemize}
  \item \textbf{Bleeding:} Intraoperative or postoperative hemorrhage, potentially requiring blood transfusion, re-operation, or interventional radiology.
  
  \item \textbf{Infection:} Surgical site infection (SSI) of the wound (superficial or deep), or deeper intra-abdominal infections such as abscess formation, which may require drainage or antibiotics.
  
  \item \textbf{Anesthesia Complications:} Adverse reactions to anesthetic agents, respiratory depression, cardiac events (e.g., myocardial infarction, arrhythmia), stroke, or pneumonia.
  
  \item \textbf{Deep Vein Thrombosis (DVT) and Pulmonary Embolism (PE):} Blood clot formation in the legs that can travel to the lungs, a potentially life-threatening complication, mitigated by prophylactic measures.
  
  \item \textbf{Organ Injury:} Accidental damage to adjacent organs such as the bladder, ureters, or major blood vessels during dissection, potentially requiring repair.
\end{itemize}

\textbf{Procedure-Specific Risks:}
\begin{itemize}
  \item \textbf{Anastomotic Leak:} The most serious complication, where the newly joined bowel segments separate, leading to leakage of bowel contents into the abdominal cavity. This can cause peritonitis, sepsis, multi-organ failure, and often requires urgent re-operation, creation of a new diverting stoma, or even permanent ostomy.
  
  \item \textbf{Small Bowel Obstruction:} Blockage of the small intestine due to adhesions formed after surgery, internal hernia, or anastomotic edema, requiring further intervention such as nasogastric tube decompression, or potentially re-operation.
  
  \item \textbf{Stoma Site Hernia (Incisional Hernia):} Weakness in the abdominal wall at the site of the previous stoma, leading to protrusion of abdominal contents, which can cause pain, obstruction, or strangulation and may require subsequent repair.
  
  \item \textbf{Wound Dehiscence:} Partial or complete separation of the surgical wound edges, often requiring wound care and sometimes re-closure.
  
  \item \textbf{Prolonged Ileus:} Delayed return of normal bowel function (more than 5-7 days), leading to nausea, vomiting, abdominal distension, and prolonged hospital stay.
  
  \item \textbf{Fecal Fistula:} An abnormal connection between the bowel and the skin or another organ, leading to persistent drainage of bowel contents, often requiring prolonged wound care and sometimes re-operation.
  
  \item \textbf{Bowel Dysfunction (Low Anterior Resection Syndrome):} Persistent changes in bowel habits, including increased frequency, urgency, clustering of bowel movements, or even fecal incontinence, particularly after low colorectal anastomoses. This can significantly impact quality of life.
  
  \item \textbf{Anastomotic Stricture:} Narrowing of the anastomotic site over time due to scarring, potentially causing obstructive symptoms and requiring endoscopic dilation or further surgery.
  
  \item \textbf{Recurrence of Underlying Disease:} While not a direct complication of reversal, the underlying condition (e.g., diverticulitis, Crohn's disease) may recur, potentially necessitating further surgical intervention or even re-creation of a stoma.
\end{itemize}

\section*{Postoperative Care \& Recovery}
Following stoma reversal, the patient is transferred to the Post-Anesthesia Care Unit (PACU) for close monitoring of vital signs, pain control, and recovery from anesthesia. Initial management focuses on adequate analgesia, intravenous fluid balance, and early mobilization to prevent complications like DVT and pneumonia. Patients are typically kept NPO initially, with a gradual advancement to clear liquids, then a soft diet, as bowel function returns, indicated by the passage of flatus or the first bowel movement per rectum. The first bowel movement is a significant milestone, confirming the patency of the new anastomosis. Most patients will experience some degree of abdominal discomfort, bloating, and changes in bowel habits in the immediate postoperative period.

Over the first 1-2 weeks post-surgery, patients are typically discharged from the hospital once they are tolerating oral intake, managing pain with oral medications, and demonstrating adequate bowel function. Wound care instructions are provided, and patients are advised to gradually increase their activity levels while avoiding heavy lifting (typically for 6-8 weeks) to protect the fascial repair. It is common for bowel habits to be irregular, with increased frequency, urgency, or even mild incontinence, especially in the initial weeks to months. This is due to the bowel adapting to its restored continuity and the absence of the diverting stoma. Long-term recovery involves the gradual normalization of bowel function, which can take several weeks to many months, and patients may require dietary adjustments, fiber supplements, or pelvic floor exercises to manage persistent symptoms. Full recovery, including return to pre-surgery activity levels, usually takes 2-3 months.

During stoma reversal, the surgeon meticulously documents the intraoperative findings, including the condition of the bowel segments (proximal and distal), the presence and extent of adhesions, the quality of the tissue at the stoma site, and the appearance of the anastomosis. Any unexpected findings, such as active inflammation, strictures, or suspicious lesions, are noted. The patency of the distal bowel is confirmed, often by digital rectal exam or intraoperative endoscopy.

If a segment of bowel, particularly the stoma-bearing portion, is resected, it is sent to pathology for histological examination. The pathology report will confirm the nature of the resected tissue, ensuring that the margins are free of inflammation or malignancy, and that the tissue is healthy. For example, it might confirm the absence of active diverticulitis or Crohn's disease in the resected segment. This report provides crucial information about the health of the tissue used for the anastomosis and helps guide future management, especially in cases where the initial stoma was created for inflammatory conditions or malignancy. The absence of significant pathological findings in the resected stoma segment is generally a favorable outcome.

A normal and successful postoperative course following stoma reversal is characterized by a predictable progression towards restored bowel function and overall well-being. Initially, patients will experience mild to moderate abdominal pain, managed effectively with oral analgesics. Bowel sounds should return within 24-48 hours, followed by the passage of flatus and then soft, unformed stools per rectum. The diet is gradually advanced from clear liquids to a regular diet as tolerated. The stoma site wound should heal cleanly, without signs of infection or dehiscence. Patients are typically discharged within 3-7 days, once they are ambulating, tolerating oral intake, and managing their pain. Over the subsequent weeks to months, bowel habits will gradually normalize, with decreasing frequency and increasing consistency of stools. While some degree of altered bowel function (e.g., increased frequency, urgency) is common, especially after low anastomoses, these symptoms typically improve significantly over 3-6 months as the bowel adapts. The ultimate expectation is the complete resolution of stoma-related issues and a significant improvement in the patient's quality of life.

Postoperative care and follow-up are essential to monitor recovery, address any complications, and support the patient's adaptation to restored bowel function.

\begin{itemize}
  \item \textbf{Postoperative Clinic Visits:} Typically scheduled for 2-4 weeks after discharge to assess wound healing, discuss bowel function, review pathology results, and address any patient concerns or symptoms. Further visits may be scheduled as needed.
  
  \item \textbf{Suture/Staple Removal:} If non-absorbable skin sutures or staples were used, they are removed during the initial follow-up visit.
  
  \item \textbf{Dietary Counseling:} Guidance on diet modification to manage bowel habits, especially in the initial months, may be provided. This can include advice on increasing fiber intake, ensuring adequate hydration, and avoiding trigger foods that exacerbate symptoms.
  
  \item \textbf{Physical Therapy/Pelvic Floor Exercises:} For patients experiencing persistent bowel dysfunction (e.g., urgency, minor incontinence), referral to a pelvic floor physical therapist may be beneficial to strengthen pelvic muscles and improve control.
  
  \item \textbf{Imaging/Endoscopy:} If complications such as anastomotic stricture, persistent leak, or recurrent obstruction are suspected, imaging studies (e.g., CT scan, contrast enema) or endoscopic evaluation (e.g., colonoscopy) may be performed.
  
  \item \textbf{Warning Signs:} Patients are educated on warning signs that necessitate immediate medical attention, including persistent fever (over 101.5\textdegree F or 38.6\textdegree C), severe or worsening abdominal pain, persistent nausea or vomiting, inability to pass gas or stool, significant wound redness, swelling, or purulent drainage, or heavy rectal bleeding.
\end{itemize}

Long-term follow-up may also involve surveillance for the underlying disease that led to the initial stoma, such as regular colonoscopy for cancer surveillance or management of inflammatory bowel disease.

\section*{Outcomes \& Prognosis}
The epidemiology of stoma reversal is intrinsically linked to the incidence of temporary stoma creation. While precise global figures are challenging to ascertain, it is estimated that hundreds of thousands of temporary stomas are created annually worldwide. Ileostomy reversals are generally more common than colostomy reversals, largely due to the frequent use of diverting loop ileostomies to protect low colorectal anastomoses following anterior resections for rectal cancer or restorative proctocolectomy for inflammatory bowel disease. 

There is no specific age or sex predilection for stoma reversal itself; rather, the demographic profile mirrors that of the underlying conditions necessitating the initial stoma, such as colorectal cancer (more common in older adults, peak incidence in 60s-70s) or inflammatory bowel disease (often diagnosed in younger populations, 15-30 years). Common indications for initial stoma creation that lead to reversal include colorectal malignancy (approximately 40-50\% of cases), complicated diverticulitis (20-30\%), inflammatory bowel disease (10-15\%), and trauma or ischemia (5-10\%). Despite being a reconstructive procedure, stoma reversal carries significant morbidity, with overall complication rates ranging from 20\% to 40\%. The most feared complication, anastomotic leak, occurs in 5-15\% of cases, depending on patient factors and the complexity of the initial surgery. Mortality rates are relatively low, typically between 1\% and 3\%, but can be higher in patients with significant comorbidities or those undergoing reversal in an emergency setting.

Reversal can restore intestinal continuity, but technical success and functional recovery are separate outcomes. Results depend on the original operation, anastomotic level, bowel and sphincter function, comorbidity, adhesions, and the reason for diversion. Patients may develop leak, obstruction, hernia, altered frequency, urgency, clustering, or incontinence; low anterior resection syndrome may persist even when the anastomosis is intact.

The prognosis for long-term survival is primarily dictated by the underlying disease that necessitated the initial stoma (e.g., stage of colorectal cancer). Recurrence of complications includes anastomotic stricture (requiring endoscopic dilation or re-operation in 5-10\% of cases) and incisional hernia at the stoma site, which can occur in 10-30\% of patients and may necessitate further surgical repair. Prognostic factors influencing both functional and surgical outcomes include patient comorbidities (e.g., diabetes, obesity, smoking), the quality of the distal bowel segment (e.g., prior radiation exposure), the type and complexity of the initial surgery, and the surgeon's experience. Despite potential functional challenges, the vast majority of patients report an improved quality of life after successful stoma reversal compared to living with a permanent ostomy, highlighting the significant benefit of this reconstructive procedure.

\section*{References}
\begin{enumerate}
  \item MedlinePlus. Colostomy reversal. U.S. National Library of Medicine; 2025. Available from: \url{https://medlineplus.gov/ency/article/002941.htm}
  
  \item Brunicardi F, Andersen DK, Billiar TR, et al. Schwartz's Principles of Surgery. 11th ed. McGraw-Hill Education; 2019.
  
  \item Sabiston DC, Townsend CM. Sabiston Textbook of Surgery: The Biological Basis of Modern Surgical Practice. 21st ed. Elsevier; 2021.
  
  \item American Society of Colon and Rectal Surgeons (ASCRS). Clinical Practice Guidelines for Ostomy Surgery. 2023. Available from: \url{https://www.fascrs.org/patients/conditions/ostomy-surgery}
\end{enumerate}

\clearpage
